\section{Expected Results}
\label{sec:results_expected}

This manuscript is currently structured around the expected outcomes of the experiments rather than final empirical values. The tables below define the result template that will be populated once the full FEMNIST pipeline is executed. Because the current codebase is designed for reproducibility, the expected outputs are already tied to fixed experimental parameters: `seed = 42`, `num_clients = 5`, `num_rounds = 3`, and `excluded_client_id = client_0` in the unlearning setting.

\subsection{Expected Training Summary}

\begin{table}[H]
\centering
\caption{Expected summary of the pure QFL training experiment.}
\label{tab:training_expected}
\begin{tabular}{lccc}
\toprule
Metric & Round 1 & Round 2 & Round 3 \\
\midrule
Number of active clients & 5 & 5 & 5 \\
Aggregation stability & High & High & High \\
Global model availability & Yes & Yes & Yes \\
GPU backend usage & Preferred & Preferred & Preferred \\
\bottomrule
\end{tabular}
\end{table}

\subsection{Expected Unlearning Summary}

\begin{table}[H]
\centering
\caption{Expected metrics for the QFI-based unlearning experiment.}
\label{tab:unlearning_expected}
\begin{tabular}{lcc}
\toprule
Metric & Target behavior & Interpretation \\
\midrule
Forget-set accuracy & Decrease toward random baseline & Removed client is effectively forgotten \\
Retain-set accuracy & Remain stable or slightly decrease & No catastrophic forgetting \\
MIA success rate & Decrease & Reduced privacy leakage \\
QFI trace & Reduce after unlearning & Lower sensitivity to removed client \\
\bottomrule
\end{tabular}
\end{table}

\subsection{Placeholder Result Schema}

The following JSON-like schema reflects the expected output generated by the unlearning script:

\begin{verbatim}
{
  "qfi_trace": <float>,
  "remaining_clients": 4.0,
  "excluded_client": "client_0",
  "forget_set_accuracy": <float>,
  "retain_set_accuracy": <float>,
  "random_baseline_accuracy": <float>,
  "mia_success_rate": <float>
}
\end{verbatim}

The final paper should replace placeholders with actual values, confidence intervals, and, if available, multiple runs across random seeds.

When the experiments are executed on the RunPod environment described in the repository, the resulting summaries should be cross-checked against the checkpoint files and logs to verify that the same deterministic configuration was used end-to-end.
